\section{The construction problem and the order of its limits}
\label{ex-statement}

The aim is to construct a relativistic field theory from finite-dimensional
Wilson integrals. There are two limiting operations. First the lattice spacing
tends to zero while the physical period stays fixed. Then the period tends to
infinity. The observables and bare action coefficients must be chosen coherently
before either operation: changing the family in response to a requested sequence
of periods would give a weaker statement.

We assume familiarity with finite-dimensional Lie groups, probability Haar
measure on compact groups, Lebesgue integration and dominated convergence,
and elementary Hilbert-space theory. The specific compactness, quotient and
operator arguments used below are supplied where they enter the construction.

Fix a closed subgroup \(G\subset U(r_G)\), with its specified faithful matrix
representation. Its real Lie algebra is
\[
 \mathfrak g=\{X\in M_{r_G}(\mathbb C):\exp(tX)\in G
                       \text{ for every }t\in\mathbb R\}.
\]
We assume that \(\mathfrak g\) is nonabelian and simple: its only real
subspaces invariant under every \(\operatorname{ad}X\), \(X\in\mathfrak g\),
are \(0\) and \(\mathfrak g\). Coordinates on \(\mathbb R^4\) are written
\(x=(\boldsymbol x,t)\), with time last. The Minkowski form is
\(t^2-|\boldsymbol x|^2\), and the closed forward cone is
\(\overline V_+=\{(\boldsymbol p,E):|\boldsymbol p|\le E\}\).

The lattice family consists of strictly increasing positive integer sequences
\(P_p\) and \(D_N\), a spacing \(a_N=D_N^{-1}\), real coefficients
\(\beta_{N,p}\), and local gauge-invariant insertion functions
\(O_{N,p,v}(x;U)\) indexed by tokens \(v\). The periodic lattice has
\(P_pD_N\) sites along each axis and physical period \(P_p\).
Section~\ref{ex-lattice} describes its Wilson probability measure and moments.
Each insertion is continuous in \(U\). For each token a single finite
integer-offset stencil controls the insertion at every spacing, period and
placement. The running condition is
\begin{equation}\label{ex-running-input}
 \exists\sigma>0\quad\forall\varepsilon>0\quad\exists N_0\quad
 \forall N\ge N_0\quad\forall p:\qquad
 \left|\frac{\beta_{N,p}}{\log(a_N^{-1})}-\sigma\right|<\varepsilon.
\end{equation}
In particular, the choice of \(N_0\) is uniform in the period index.

Let \(\operatorname{Conf}_n(\mathbb R^4)\) be the set of ordered
\(n\)-tuples of pairwise distinct points. A separated test is a smooth
complex function with compact support in this set, extended by zero outside
it. For insertion tokens \(\boldsymbol v=(v_1,\ldots,v_n)\), write
\(M_{N,p}(\boldsymbol v;\boldsymbol x)\) for their Wilson moment and put
\begin{equation}\label{ex-smearing}
 W_{N,p}(\boldsymbol v;f)=a_N^{4n}
 \sum_{\boldsymbol x\in\Lambda_{N,p}^n}
 f(a_N\boldsymbol x)M_{N,p}(\boldsymbol v;\boldsymbol x).
\end{equation}
Here the coordinates of \(\boldsymbol x\) are chosen in the centered
fundamental cell. A period is suitable for \(f\) if its support lies strictly
inside that physical cell in every coordinate.

The limiting data have several related components. For every integer grade
\(d\ge0\) there is a finite-dimensional complex space \(\mathcal F_d\).
It has a complex-linear injection into a common field-label space; every label is
represented at some grade. Tokens map onto the disjoint union
\(\coprod_d\mathcal F_d\). A compact metric space \(X\), a probability
measure \(\nu\) of full support, and a selected character \(\chi_*\in X\)
describe the phases. For graded fields \(\boldsymbol A\), let
\(S^{\chi}_{\boldsymbol A}(\boldsymbol x)\) and
\(S^{\mathrm{mix}}_{\boldsymbol A}(\boldsymbol x)\) denote the phase and
mixed separated correlations. They are smooth on configuration space,
permutation symmetric, and continuous in \(\chi\), with
\begin{equation}\label{ex-mixture}
 S^{\mathrm{mix}}_{\boldsymbol A}(\boldsymbol x)
   =\int_X S^{\chi}_{\boldsymbol A}(\boldsymbol x)\,d\nu(\chi).
\end{equation}
The mixed correlations are the Wilson limits. The selected correlations
are the Euclidean correlations to be continued into the reconstructed theory.
These two roles should be kept distinct even when the phase space has only
one point.

\begin{definition}[Reconstructed selected theory]\label{ex-reconstructed}
The selected theory consists of a complex Hilbert space \(\mathcal H\),
a normalized vacuum \(\Omega\), a common dense invariant domain
\(\mathcal D\) containing \(\Omega\), and complex-linear fields
\(\Phi(A,f):\mathcal D\to\mathcal D\), linear separately in the label
and test, where
\(f\in\mathcal S(\mathbb R^4)\). For fixed \(A,\psi\), the map
\(f\mapsto\Phi(A,f)\psi\) is continuous in the Schwartz topology.
The domain is the span of field words applied to the vacuum. There is a
conjugate-linear label adjoint \(A\mapsto A^*\) with
\[
 \langle\Phi(A,f)\psi,\eta\rangle
   =\langle\psi,\Phi(A^*,\overline f)\eta\rangle.
\]
A strongly continuous unitary representation of the proper orthochronous
Poincar\'e group preserves the domain and the vacuum, and transforms fields
with the specified action on labels and the ordinary affine pullback of tests.
Its translation spectral measure is supported in \(\overline V_+\), and
every translation-invariant vector is a scalar multiple of \(\Omega\).
Fields commute on \(\mathcal D\) when their test supports are spacelike
separated. Every nonzero label acts nontrivially for some test, and at least
one smeared field is nonzero.

For fixed spatial positions the selected correlations have a holomorphic
continuation \(F_{\boldsymbol A}(\boldsymbol z;\boldsymbol x)\) on
\(\operatorname{Im}z_1<\cdots<\operatorname{Im}z_n\). At strictly ordered
Euclidean times its value at \(z_i=it_i\) is
\(S^{\chi_*}_{\boldsymbol A}((\boldsymbol x_i,t_i)_i)\).
For every strictly increasing real height vector \(\boldsymbol h\) and
every tuple of Schwartz tests, its integrals on
\(z_i=t_i+i\epsilon h_i\) converge as \(\epsilon\downarrow0\) to the
corresponding vacuum expectation of the smeared fields.
\end{definition}

The existence target is one regular Wilson family such that, for every
strictly increasing period-index sequence \(p(k)\), a strictly increasing
\(q(k)\) and one hierarchy with the preceding properties can be found.
For each separated test there are row limits \(w_p(\boldsymbol v;f)\)
satisfying
\begin{align}
 \lim_{N\to\infty}W_{N,p}(\boldsymbol v;f)
   &=w_p(\boldsymbol v;f)
       &&\text{at each suitable fixed }p,\label{ex-row-target}\\
 \lim_{k\to\infty}w_{p(q(k))}(\boldsymbol v;f)
   &=\int f(\boldsymbol x)
           S^{\mathrm{mix}}_{\boldsymbol A}(\boldsymbol x)\,d\boldsymbol x,
       \label{ex-volume-target}
\end{align}
where \(A_i\) is the class of \(v_i\). Suitability holds eventually
along the selected period sequence. The same \(q\) and hierarchy work for
all arities, tokens and tests. The hierarchy may depend on the requested
sequence; the Wilson family may not.

The full endpoint also specifies a common operator-product expansion and a
short-distance trace test, described in Section~\ref{ex-assembly}. These
are additional obligations on the same hierarchy. The target concerns a
theory of the constructed gauge-invariant observables; identifying a particular
label with a separately prescribed physical observable requires its own
construction argument.
